\begin{pblproject}
\label{sec:12-path-algebras:07-checkpoint-project:1}
\textbf{\large Checkpoint project: path length, and composition respects it}\par\smallskip


Part II closes here, having built groups, rings, modules, and now quivers and paths from scratch. This project extends the \texttt{Path}/\texttt{Path.append} construction from Sections 4--5 with one more piece of structure, the \textbf{length} of a path, and proves it behaves the way it obviously should under composition, tying together the inductive-type work of this chapter with the ``prove it once, generically'' habit Chapters 8 and 10 established.

\textbf{Learning objectives.} Define a recursive function over the same indexed inductive type \texttt{Path} (Section 4), then prove a genuine theorem about it by induction, mirroring the recursion of \texttt{Path.append} itself case for case (Section 5). Along the way, this project surfaces a real, verifiable fact about Lean, because \texttt{Path} is \emph{indexed} by both its endpoints, functions matching on it (both \texttt{Path.append} and the \texttt{Path.length} built here) do not reduce by bare \texttt{rfl} once an abstract path variable is involved. Only their auto-generated equation lemmas do. This is worth discovering directly rather than being told, the same way the worked proofs of Section 5 are best understood by predicting what each \texttt{rw} will do before running it.

\textbf{Prerequisites.} Chapters 1--12, specifically Section 4 (\texttt{Path}) and Section 5 (\texttt{Path.append}).

\textbf{Milestones.}

\begin{enumerate}
\def\labelenumi{\arabic{enumi}.}
\tightlist
\item
  Define \texttt{Path.length : \{u v : V\} → Path Q u v → Nat} by recursion, the trivial path \texttt{nil} has length \texttt{0}; \texttt{cons a h h' p} has length \texttt{p.length + 1} (one more than the path it extends).
\item
  Compute a couple of lengths by hand and check them with \texttt{\#eval} against the example paths of Section 4 (\texttt{pathAlpha}, \texttt{pathBetaAlpha}).
\item
  Prove \texttt{Path.append\_\allowbreak{}length : (Path.append p q).length = p.length + q.length} by induction on \texttt{q}, mirroring the \texttt{nil}/\texttt{cons} cases of \texttt{Path.append} itself. Predict, before running it, whether each case closes by \texttt{rfl}, then discover (as the note above flags) that it does not, and that \texttt{simp only [Path.append, Path.length]} is needed to unfold one step, exactly as \texttt{rfl} needs the \emph{equation lemmas} Lean generates for an indexed match, not raw iota-reduction, once an abstract path is involved.
\end{enumerate}

\textbf{Deliverable.} \texttt{Path.length} and the proved theorem \texttt{Path.append\_\allowbreak{}length}, checked against at least one concrete instance.

\textbf{Self-verification.}

\begin{lstlisting}
def Path.length {V A : Type} {Q : Quiver V A} : {u v : V} (*$\rightarrow$*) Path Q u v (*$\rightarrow$*) Nat
  | _, _, Path.nil _ => 0
  | _, _, Path.cons _ _ _ p => p.length + 1

#eval pathAlpha.length        -- 1
#eval pathBetaAlpha.length    -- 2

theorem Path.append_length {V A : Type} {Q : Quiver V A} {u v w : V}
    (p : Path Q u v) (q : Path Q v w) :
    (Path.append p q).length = p.length + q.length := by
  induction q with
  | nil =>
    simp only [Path.append, Path.length]
    rw [Nat.add_zero]
  | cons a h h' q' ih =>
    simp only [Path.append, Path.length]
    rw [ih, Nat.add_assoc]

-- The concrete check: pathBetaAlpha was Section 5's own worked example of
-- Path.append; its length should be pathAlpha.length + pathBetaOnly.length.
example : (Path.append pathAlpha pathBetaOnly).length =
    pathAlpha.length + pathBetaOnly.length :=
  Path.append_length pathAlpha pathBetaOnly

#eval (Path.append pathAlpha pathBetaOnly).length   -- 2
\end{lstlisting}

\texttt{Path.append} and \texttt{Path.length} both recurse, and neither result can be printed on its own. A bare \texttt{Path} has no \texttt{Repr} instance (there is no generic way to display the arrows of an arbitrary quiver as text), which is exactly why every example above prints a \texttt{.length} rather than the path itself. Composing the two together, with a \texttt{dbg\_\allowbreak{}trace} in each, makes both recursions visible in the one place they can actually be observed.

\begin{lstlisting}
def Path.append' {V A : Type} {Q : Quiver V A} {u v w : V}
    (p : Path Q u v) (q : Path Q v w) : Path Q u w :=
  match q with
  | Path.nil _ => dbg_trace s!"append: q is nil, base case, returning p unchanged"; p
  | Path.cons a h h' q' =>
      dbg_trace s!"append: q ends with an arrow, recursing on the shorter path underneath, then re-attaching that arrow";
      Path.cons a h h' (Path.append' p q')

def Path.length' {V A : Type} {Q : Quiver V A} : {u v : V} (*$\rightarrow$*) Path Q u v (*$\rightarrow$*) Nat
  | _, _, Path.nil _ => dbg_trace s!"length: nil, base case, returning 0"; 0
  | _, _, Path.cons _ _ _ p => dbg_trace s!"length: cons, adding 1 to the rest's length"; p.length' + 1

#eval (Path.append' pathAlpha pathBetaOnly).length'
-- append: q ends with an arrow, recursing on the shorter path underneath, then re-attaching that arrow
-- append: q is nil, base case, returning p unchanged
-- length: cons, adding 1 to the rest's length
-- length: cons, adding 1 to the rest's length
-- length: nil, base case, returning 0
-- 2
\end{lstlisting}

Read the five lines in order. The first two are \texttt{Path.append'} building the composed path (it recurses on \texttt{q = pathBetaOnly}, prints once for its one \texttt{cons}, then once for the trailing \texttt{nil}, matching the length of \texttt{pathBetaOnly} itself, which is 1). Only once that path exists do the next three lines run \texttt{Path.length'} over it, one \texttt{cons}/\texttt{cons}/\texttt{nil} line per constructor of the freshly-built two-arrow path, arriving at the same \texttt{2} the untraced version above already computed. \texttt{Path.append} finishes building the whole structure \emph{before} \texttt{Path.length} ever starts walking it. The two recursions do not interleave step-by-step, even though both traces appear in one \texttt{\#eval}.

If this compiles and the \texttt{\#eval}s match, the project is done. A full worked solution is in \hyperref[sec:15-appendix-solutions:12-chapter-12:1]{Appendix, Chapter 12}.

\end{pblproject}
